% =============================================================================
% preamble.tex — Pebble Whitepaper
% =============================================================================

% ---------------------------------------------------------------------------
% Core packages
% ---------------------------------------------------------------------------
\usepackage[T1]{fontenc}
\usepackage[utf8]{inputenc}
\usepackage{lmodern}
\usepackage{microtype}
\usepackage{geometry}
\geometry{a4paper, margin=1in}

% ---------------------------------------------------------------------------
% Math
% ---------------------------------------------------------------------------
\usepackage{amsmath}
\usepackage{amssymb}
\usepackage{amsthm}
\usepackage{mathtools}

% ---------------------------------------------------------------------------
% Graphics and diagrams
% ---------------------------------------------------------------------------
\usepackage{tikz}
\usepackage{tikz-cd}
\usetikzlibrary{
  arrows.meta,
  shapes.geometric,
  positioning,
  fit,
  backgrounds,
  calc,
  decorations.pathreplacing,
  matrix
}
\usepackage{pgfplots}
\pgfplotsset{compat=1.18}
\usepackage{graphicx}
\usepackage{float}

% ---------------------------------------------------------------------------
% Tables and listings
% ---------------------------------------------------------------------------
\usepackage{booktabs}
\usepackage{array}
\usepackage{longtable}
\usepackage{listings}
\usepackage{xcolor}

% ---------------------------------------------------------------------------
% References and links
% ---------------------------------------------------------------------------
\usepackage{hyperref}
\hypersetup{
  colorlinks=true,
  linkcolor=blue!70!black,
  citecolor=green!50!black,
  urlcolor=blue!70!black,
  pdfauthor={Abderrahmane Memmou},
  pdftitle={Pebble: A Leaderless, Eventually Consistent Distributed Key-Value Store},
  pdfsubject={Distributed Systems},
}
\usepackage[nameinlink]{cleveref}

% ---------------------------------------------------------------------------
% Algorithms
% ---------------------------------------------------------------------------
\usepackage{algorithm}
\usepackage{algpseudocode}

% ---------------------------------------------------------------------------
% Typography
% ---------------------------------------------------------------------------
\usepackage{parskip}
\setlength{\parskip}{6pt}
\usepackage{enumitem}
\usepackage{caption}
\usepackage{subcaption}

% ---------------------------------------------------------------------------
% Code listings style
% ---------------------------------------------------------------------------
\definecolor{codebg}{RGB}{248,248,248}
\definecolor{codeframe}{RGB}{200,200,200}
\definecolor{codecomment}{RGB}{100,100,100}
\definecolor{codekeyword}{RGB}{0,0,180}
\definecolor{codestring}{RGB}{163,21,21}
\definecolor{codenumber}{RGB}{100,100,100}

\lstdefinestyle{pebblestyle}{
  backgroundcolor=\color{codebg},
  frame=single,
  rulecolor=\color{codeframe},
  basicstyle=\ttfamily\small,
  keywordstyle=\color{codekeyword}\bfseries,
  stringstyle=\color{codestring},
  commentstyle=\color{codecomment}\itshape,
  numberstyle=\tiny\color{codenumber},
  numbers=left,
  numbersep=8pt,
  breaklines=true,
  captionpos=b,
  tabsize=2,
  showstringspaces=false,
}
\lstset{style=pebblestyle}

\lstdefinelanguage{json}{
  morestring=[b]",
  morecomment=[l]{//},
  morecomment=[s]{/*}{*/},
  literate=
    *{0}{{{\color{codenumber}0}}}{1}
     {1}{{{\color{codenumber}1}}}{1}
     {2}{{{\color{codenumber}2}}}{1}
     {3}{{{\color{codenumber}3}}}{1}
     {4}{{{\color{codenumber}4}}}{1}
     {5}{{{\color{codenumber}5}}}{1}
     {6}{{{\color{codenumber}6}}}{1}
     {7}{{{\color{codenumber}7}}}{1}
     {8}{{{\color{codenumber}8}}}{1}
     {9}{{{\color{codenumber}9}}}{1}
     {:}{{{\color{codekeyword}{:}}}}{1}
     {,}{{{\color{codekeyword}{,}}}}{1}
     {\{}{{{\color{codekeyword}{\{}}}}{1}
     {\}}{{{\color{codekeyword}{\}}}}}{1}
     {[}{{{\color{codekeyword}{[}}}}{1}
     {]}{{{\color{codekeyword}{]}}}}{1},
}

% ---------------------------------------------------------------------------
% Theorem environments
% ---------------------------------------------------------------------------
\theoremstyle{plain}
\newtheorem{theorem}{Theorem}[section]
\newtheorem{lemma}[theorem]{Lemma}
\newtheorem{corollary}[theorem]{Corollary}
\newtheorem{proposition}[theorem]{Proposition}

\theoremstyle{definition}
\newtheorem{definition}[theorem]{Definition}
\newtheorem{example}[theorem]{Example}
\newtheorem{invariant}[theorem]{Invariant}
\newtheorem{protocol}[theorem]{Protocol}

\theoremstyle{remark}
\newtheorem{remark}[theorem]{Remark}
\newtheorem{note}[theorem]{Note}

% ---------------------------------------------------------------------------
% Custom macros
% ---------------------------------------------------------------------------
% Mutation fields
\newcommand{\mutid}{\mathrm{id}}
\newcommand{\mnid}{\mathit{nodeId}}
\newcommand{\mseq}{\mathit{seq}}
\newcommand{\mts}{\mathit{ts}}
\newcommand{\mop}{\mathit{op}}
\newcommand{\mkey}{\mathit{key}}
\newcommand{\mval}{\mathit{val}}
\newcommand{\msig}{\sigma}

% Sets and domains
\newcommand{\NodeSet}{\mathcal{N}}
\newcommand{\MutSet}{\mathcal{M}}
\newcommand{\KeySet}{\mathcal{K}}
\newcommand{\VV}{\mathbf{V}}

% LWW order
\newcommand{\lwwlt}{\prec_{\mathrm{lww}}}
\newcommand{\lwwleq}{\preceq_{\mathrm{lww}}}
\newcommand{\win}{\mathrm{winner}}

% Complexity
\newcommand{\Oof}[1]{O\!\left(#1\right)}

% Protocol names
\newcommand{\pebble}{\textsc{Pebble}}
\newcommand{\pkg}[1]{\texttt{@pebbl/#1}}

% Inline code
\newcommand{\code}[1]{\texttt{#1}}

% Message types
\newcommand{\msg}[1]{\textsf{\small #1}}

% TikZ styles
\tikzset{
  node box/.style={
    draw, rectangle, rounded corners=4pt,
    minimum width=3cm, minimum height=0.8cm,
    fill=blue!8, draw=blue!40, thick,
    font=\small
  },
  msg arrow/.style={
    -Stealth, thick, draw=black!70
  },
  dashed arrow/.style={
    -Stealth, dashed, draw=black!50
  },
  component/.style={
    draw, rectangle, rounded corners=6pt,
    minimum width=4.5cm, minimum height=1cm,
    fill=gray!10, draw=gray!60, thick,
    align=center,
    font=\small\bfseries
  },
  layer/.style={
    draw, rectangle, rounded corners=3pt,
    fill=orange!8, draw=orange!40, thick,
    font=\footnotesize
  },
}
